\documentclass[12pt,oneside]{book}

% ==========================================================
% METADATOS
% ==========================================================
\newcommand{\universidad}{Universidad de Buenos Aires}
\newcommand{\facultad}{Facultad de Derecho}
\newcommand{\carrera}{Consultorio Jurídico Gratuito --- Área de Familia}
\newcommand{\materia}{Práctica Profesional}
\newcommand{\materiacompleta}{Práctica Profesional --- Patrocinio Jurídico Gratuito, Derecho de Familia}
\newcommand{\titulo}{Patrocinio Jurídico Gratuito --- Derecho de Familia}
\newcommand{\autor}{Isaac Edgar Camacho Ocampo}
\newcommand{\anio}{2026}

% ==========================================================
% PAQUETES
% ==========================================================
\usepackage[utf8]{inputenc}
\usepackage[T1]{fontenc}
\usepackage[spanish,es-nodecimaldot]{babel}

\usepackage[
    top=2.5cm,
    bottom=2.5cm,
    left=3cm,
    right=2.5cm,
    headheight=15pt
]{geometry}

\usepackage{amsmath}
\usepackage{amssymb}
\usepackage{mathtools}

\usepackage{graphicx}
\usepackage{float}
\usepackage{caption}
\usepackage{subcaption}

\usepackage{booktabs}
\usepackage{array}
\usepackage{longtable}
\usepackage{tabularx}

\usepackage{enumitem}

\usepackage{xcolor}
\usepackage[most]{tcolorbox}

\usepackage{fancyhdr}
\usepackage{ragged2e}

\usepackage[
    colorlinks=true,
    linkcolor=blue!50!black,
    citecolor=green!40!black,
    urlcolor=blue!60!black,
    pdftitle={\titulo},
    pdfauthor={\autor},
    pdfsubject={Apuntes universitarios},
    bookmarks=true
]{hyperref}

% ==========================================================
% CONFIGURACIÓN
% ==========================================================
\graphicspath{
    {./img/}
    {./graficos/}
    {./figuras/}
}

\setcounter{tocdepth}{2}
\setcounter{secnumdepth}{0}

\setlength{\parindent}{0pt}
\setlength{\parskip}{0.5em}

\setlist[itemize]{
    topsep=0.3em,
    itemsep=0.2em
}

\setlist[enumerate]{
    topsep=0.3em,
    itemsep=0.2em
}

\clubpenalty=10000
\widowpenalty=10000

% ==========================================================
% COMANDOS
% ==========================================================
\newcommand{\abs}[1]{\left|#1\right|}
\newcommand{\norm}[1]{\left\lVert#1\right\rVert}
\newcommand{\vect}[1]{\mathbf{#1}}
\newcommand{\deriv}[2]{\frac{d#1}{d#2}}
\newcommand{\pderiv}[2]{\frac{\partial#1}{\partial#2}}

% ==========================================================
% ENTORNOS / CAJAS
% ==========================================================
\newtcolorbox{definition}[1]{
    enhanced,
    breakable,
    title={Definición: #1},
    fonttitle=\bfseries,
    colback=blue!3,
    colframe=blue!50!black,
    boxrule=0.8pt,
    arc=2mm
}

\newtcolorbox{important}{
    enhanced,
    breakable,
    title={Importante},
    fonttitle=\bfseries,
    colback=yellow!8,
    colframe=orange!70!black,
    boxrule=0.8pt,
    arc=2mm
}

\newtcolorbox{example}[1]{
    enhanced,
    breakable,
    title={Ejemplo: #1},
    fonttitle=\bfseries,
    colback=green!4,
    colframe=green!40!black,
    boxrule=0.8pt,
    arc=2mm
}

\newtcolorbox{formula}[1]{
    enhanced,
    breakable,
    title={Fórmula / Regla: #1},
    fonttitle=\bfseries,
    colback=purple!4,
    colframe=purple!50!black,
    boxrule=0.8pt,
    arc=2mm
}

\newtcolorbox{exercise}[1]{
    enhanced,
    breakable,
    title={Ejercicio: #1},
    fonttitle=\bfseries,
    colback=gray!5,
    colframe=gray!60!black,
    boxrule=0.8pt,
    arc=2mm
}

\newtcolorbox{note}{
    enhanced,
    breakable,
    title={Nota},
    fonttitle=\bfseries,
    colback=white,
    colframe=gray!50,
    boxrule=0.6pt,
    arc=2mm
}

% ==========================================================
% CONFIGURACIÓN FINAL
% ==========================================================
\pagestyle{fancy}
\fancyhf{}
\fancyhead[L]{\small\nouppercase{\leftmark}}
\fancyhead[R]{\small\materia}
\fancyfoot[C]{\thepage}
\renewcommand{\headrulewidth}{0.4pt}
\renewcommand{\footrulewidth}{0pt}

\fancypagestyle{plain}{
    \fancyhf{}
    \fancyfoot[C]{\thepage}
    \renewcommand{\headrulewidth}{0pt}
}

\renewcommand{\arraystretch}{1.15}

% ==========================================================
\begin{document}
% ==========================================================

\frontmatter

% ---------------------------- PORTADA ----------------------------
\begin{titlepage}

\thispagestyle{empty}

\begin{center}

\vspace*{1cm}

{\large \textsc{\universidad}}

\vspace{0.2cm}

{\large \textsc{\facultad}}

\vspace{0.2cm}

{\large \carrera}

\vspace{3cm}

{\Huge \textbf{\materia}}

\vspace{1cm}

{\LARGE \titulo}

\vspace{0.8cm}

\textit{Cada expediente estudiado con rigor es una oportunidad para aprender a decidir con criterio.}

\vspace{1.2cm}

\rule{0.70\textwidth}{0.5pt}

\vspace{0.5cm}

{\large Apuntes de estudio}

\vspace{0.5cm}

\rule{0.70\textwidth}{0.5pt}

\vfill

{\large \autor}

\vspace{0.3cm}

{\large \anio}

\end{center}

\vfill

\begin{flushleft}
\textit{Este es un modesto aporte a los estudiantes de ``\materiacompleta'' no pretende sustituir el material oficial de la materia.}
\end{flushleft}

\end{titlepage}

% ---------------------------- ÍNDICE ----------------------------
\tableofcontents

\mainmatter

% ==========================================================
% PRESENTACIÓN
% ==========================================================
\chapter*{Presentación}
\addcontentsline{toc}{chapter}{Presentación}
\markboth{Presentación}{Presentación}

Este documento reúne el contenido académico de una reunión de seguimiento y supervisión de expedientes del equipo de patrocinio jurídico gratuito del Área de Familia del Consultorio Jurídico Gratuito de la Facultad de Derecho de la Universidad de Buenos Aires (ciclo lectivo 2026). En ella se repasó el estado procesal de catorce expedientes reales de derecho de familia, agrupados en cuatro ejes:

\begin{enumerate}
    \item la constitución de la obligación alimentaria y la prueba del caudal económico;
    \item el incumplimiento y las vías de ejecución y coerción frente al deudor alimentario;
    \item el régimen de comunicación, la violencia familiar y la protección de niñas, niños y adolescentes;
    \item la filiación internacional, la capacidad jurídica y los amparos de salud vinculados a la discapacidad.
\end{enumerate}

Se identifican y explican institutos como los alimentos provisorios y definitivos, la responsabilidad solidaria entre parientes, la prescripción de la deuda alimentaria, la retención directa de haberes, el embargo, las astreintes, el exhorto internacional, la información sumaria, la revisión periódica de la sentencia de determinación de capacidad y el rol del Ministerio Público Tutelar.

\section*{Hilo conductor de los cuatro bloques}

Los catorce expedientes no son casos aislados: constituyen distintas etapas de una misma cadena de sostén económico y afectivo que el derecho de familia intenta ordenar cuando el vínculo entre los adultos se rompe. Si una casa se queda sin uno de sus sostenes, primero debe asegurarse que alguien pague las cuentas (Capítulo~1, alimentos); si esa persona no paga, hay que encontrar la forma de obligarla o de cobrar por otra vía (Capítulo~2, incumplimiento y ejecución); mientras tanto, hay que decidir quién ve a los niños, cómo y en qué condiciones, sobre todo si hubo violencia (Capítulo~3, régimen de comunicación y violencia familiar); y, si además alguna de las partes vive en otro país, tiene una discapacidad o no puede ejercer plenamente su capacidad jurídica, aparecen herramientas procesales especiales (Capítulo~4, filiación internacional, capacidad y amparo de salud). Los cuatro bloques comparten un mismo eje: la protección efectiva de quien depende económica o afectivamente de otro.

\section*{Utilidad del material para la práctica profesional}

El material no es una exposición teórica, sino el registro de cómo un equipo en formación toma decisiones estratégicas reales, bajo incertidumbre, con expedientes que no siempre responden a lo que dicen los manuales. Observar cómo se prueba el caudal económico de un deudor renuente, cómo se articula un exhorto cuando la contraparte vive en el extranjero o cuándo conviene frenar un reclamo en lugar de impulsarlo enseña el razonamiento estratégico del litigio de familia. Funciona, además, como un banco de precedentes de gestión frente a problemas recurrentes (domicilios desconocidos, deudores aparentemente insolventes, organismos de protección). Y deja una enseñanza que trasciende lo jurídico: los conflictos familiares rara vez tienen una sola versión de los hechos, y sostener una estrategia razonada y honesta es una habilidad aplicable a cualquier ámbito con personas en conflicto.

\begin{note}
La reunión consistió en un repaso expediente por expediente, a cargo de las y los estudiantes responsables de cada causa, bajo la supervisión de los abogados a cargo. Los apellidos y nombres que aparecen como título de cada unidad corresponden a la identificación interna de cada expediente.
\end{note}

% ==========================================================
% CAPÍTULO 1
% ==========================================================
\chapter[Alimentos: constitución y prueba del caudal]{Alimentos: constitución de la obligación y prueba del caudal económico}

Este primer bloque reúne los expedientes en los que la discusión gira en torno a dos preguntas iniciales de todo reclamo alimentario: quién está obligado a pagar y cómo se demuestra cuánto puede pagar. Se repasan un caso de responsabilidad solidaria familiar, un caso de alimentos subsidiarios contra un hermano ante la imposibilidad de ubicar al padre, dos casos de prueba del caudal económico mediante oficios a organismos públicos y privados, y un caso de filiación que deriva en un reclamo alimentario contra los abuelos paternos.

% ----------------------------------------------------------
\section{Unidad 1 --- Sesión inaugural: sucesión, alimentos definitivos y responsabilidad solidaria de parientes}

\subsection{Antecedentes}

La reunión se abrió con una consulta sobre la apertura de una sucesión. La abuela de la consultante falleció en la provincia de Buenos Aires y la consultante tiene su domicilio real también en la provincia, por lo que el equipo de patrocinio no podía tramitarla: el patrocinio de la Facultad no lleva causas fuera de la Ciudad Autónoma de Buenos Aires cuando el domicilio real de la parte está radicado en el conurbano o en el interior. La sucesión debía iniciarse, entonces, en el fuero provincial correspondiente.

El primer expediente propiamente dicho de la agenda fue el de una joven beneficiaria de alimentos definitivos. Cuando se habla de ``los definitivos'' sin más aclaración, se hace referencia a los que en la demanda se caratulan simplemente como ``alimentos'', por oposición a los alimentos provisorios.

\begin{definition}{Alimentos provisorios y alimentos definitivos}
Los \textbf{alimentos provisorios} se tramitan como medida cautelar dentro de la misma causa o de una conexa. Los \textbf{alimentos definitivos} son los que se reclaman en la demanda caratulada simplemente como ``alimentos''.
\end{definition}

\subsection{Situación planteada}

Falleció la abuela con quien convivía la beneficiaria. En ese marco, el equipo solicitó que los alimentos, que hasta entonces se depositaban en una cuenta a nombre de la abuela, se transfirieran a una cuenta a nombre de la propia alimentada. Al tramitarse ese pedido, el empleador del alimentante informó que este ya no trabajaba con él en relación de dependencia. Consultado el alimentante, manifestó que se había jubilado a través de la Caja de la Policía de la Provincia de Buenos Aires y que, por ser la beneficiaria mayor de veintiún años, ya no le correspondía seguir abonando alimentos, salvo que ella acreditara estar cursando estudios.

Frente a ese planteo, el equipo presentó comprobantes de que la joven realiza cursos de actualización profesional y desarrolló un argumento adicional: la beneficiaria padece convulsiones y una situación de discapacidad que, según se sostuvo, debería hacer prevalecer un deber de \textbf{responsabilidad solidaria} entre parientes, dado que además de estudiar convive con una hermana con discapacidad y colabora en su cuidado.

\subsection{Debate: probabilidades de éxito}

Un compañero adelantó que probablemente el pedido sería rechazado. Señaló que la carrera de nivel terciario que la joven ya posee (identificada en la reunión, tras una corrección conjunta, como la de asistente terapéutica) le permitiría trabajar, y que el argumento de que continúa formándose mediante cursos de especialización probablemente no alcance para sostener la obligación alimentaria más allá de los veintiún años frente a un juzgado que puede razonar que, título en mano, corresponde que la persona se sostenga por sus propios medios.

Sobre la vía de la responsabilidad solidaria entre parientes se preguntó si existían antecedentes jurisprudenciales que la admitieran con fundamento en una situación de vulnerabilidad extrema (estudios en curso más una persona con discapacidad a cargo). Quien exponía reconoció no haber encontrado fallos en ese sentido y anticipó que, a su criterio, también sería rechazada, entre otras razones porque, según lo relevado, para reclamar contra un pariente distinto del progenitor directo correspondería iniciar un incidente en el expediente principal, y ese incidente no podría ser tramitado por el equipo si la parte actora no tiene domicilio real en la Ciudad.

\begin{important}
\textbf{Responsabilidad solidaria entre parientes en materia alimentaria.} La obligación alimentaria recae en primer término sobre los progenitores; solo de manera subsidiaria y ante la imposibilidad de estos puede extenderse a otros parientes por consanguinidad (abuelos, hermanos). Invocar una responsabilidad ``solidaria'' entre parientes, como se plantea en este caso a partir de la discapacidad de la beneficiaria, es una estrategia argumental poco consolidada en la jurisprudencia relevada hasta el momento del análisis.
\end{important}

Se consultó además si existía una causa penal vinculada a la beneficiaria. Se aclaró que dicha causa no involucra al equipo ni a la parte alimentante, sino que corresponde a un trámite de guarda que la propia beneficiaria gestiona en la provincia de Buenos Aires respecto de su hermana con discapacidad, lo que podría eventualmente habilitarla a solicitar la Asignación Universal por Hijo, aunque esta no alcanzaría por sí sola para sostener los gastos del hogar.

\begin{note}
\textbf{Enseñanza para la práctica profesional.} Aun cuando la parte alimentada atraviese una situación humana especialmente difícil, la abogacía de familia exige distinguir con precisión qué puede sostenerse jurídicamente y qué corresponde a otro fuero, a otra jurisdicción o a otro tipo de asistencia estatal. La honestidad frente al pronóstico, incluso cuando resulta dolorosa para la parte, forma parte del ejercicio profesional responsable.
\end{note}

% ----------------------------------------------------------
\section{Unidad 2 --- Carlovisa: discapacidad, deudor radicado en Paraguay y alimentos contra el hermano}

\subsection{Antecedentes}

El expediente se presentó como un caso en desarrollo, con la demanda de alimentos aún en etapa de corrección: se había pedido a la estudiante responsable que corrigiera el escrito sin que hasta ese momento se hubiera recibido la versión corregida, lo que motivó un breve intercambio sobre la importancia del seguimiento activo de las correcciones solicitadas.

\subsection{Situación de la niña y contexto familiar}

El padre de la niña beneficiaria se trasladó de un día para el otro a Paraguay, donde, según se informó, reside con su hermana y con su madre. Se indicó que atravesaría además una problemática de consumo (se mencionó específicamente el alcohol). La niña presenta una discapacidad relevante, con retraso madurativo, problemas de columna y de piel, y dificultades de aprendizaje.

\begin{important}
\textbf{Certificado Único de Discapacidad (CUD) y modelo social de la discapacidad.} Para fundar el reclamo no es imprescindible contar con el CUD: basta con acreditar la discapacidad por otros medios, en línea con el modelo social de la discapacidad, que pone el eje en la existencia de barreras que limitan la participación plena de la persona y no en un documento administrativo determinado. En el caso, la niña cuenta además con la pensión por discapacidad correspondiente.
\end{important}

\subsection{Estrategia: alimentos subsidiarios contra el hermano del progenitor}

No fue posible notificar ni ubicar al padre en Paraguay (en otro expediente similar, el de Gamarra, tampoco se logró que las autoridades paraguayas informaran un domicilio pese a un exhorto internacional). Por ello el equipo decidió no iniciar una nueva causa contra el progenitor, para evitar quedar trabado durante meses a la espera de un registro civil extranjero, y promovió en cambio un reclamo de alimentos provisorios contra el hermano del padre, domiciliado en la Argentina.

Reclamar alimentos contra parientes distintos de los progenitores, en este caso un tío, es una vía ``bastante excepcional'': el orden de prelación exige intentar primero contra los abuelos y, solo ante la imposibilidad de dar con ellos (ambos abuelos residen o habrían residido en Paraguay y uno de ellos estaría fallecido), corresponde avanzar contra otros parientes colaterales, como el hermano del alimentante.

Del hermano solo se cuenta con la información obtenida de un informe de Nosis y de su perfil de Facebook, y la demanda debió estimar de forma aproximada los montos de alquiler y medicamentos, a falta de documentación más precisa, para que sea el juzgado quien determine finalmente la cuota.

\begin{important}
\textbf{Los dos extremos que exige acreditar un reclamo alimentario.} Todo reclamo de alimentos exige probar, por un lado, el caudal económico de quien debería pagar (cuánto puede pagar) y, por otro, la necesidad concreta de quien los reclama (cuánto necesita). En este caso la dificultad estuvo en ambos frentes: no hay información fehaciente sobre los ingresos del hermano demandado y los gastos de la niña debieron estimarse de manera aproximada.
\end{important}

Quedó pendiente resolver cómo acreditar el vínculo de parentesco entre el hermano demandado y el padre de la niña, ya que no se cuenta con las partidas de nacimiento paraguayas ni es sencillo obtenerlas desde la Argentina; se evaluó ofrecer prueba testimonial adicional sobre ese vínculo.

\begin{note}
\textbf{Enseñanza para la práctica profesional.} Cuando el deudor alimentario reside en el exterior y las autoridades de ese país no colaboran, la estrategia procesal debe reorientarse hacia quienes sí están al alcance de la jurisdicción argentina, incluso si ello implica invocar figuras subsidiarias y menos habituales, como la responsabilidad de un pariente colateral.
\end{note}

% ----------------------------------------------------------
\section{Unidad 3 --- Sela y García: prueba del caudal económico mediante oficios}

\subsection{Dos expedientes en espejo}

Se presentaron en conjunto dos causas de estructura muy similar, correspondientes a las consultantes Sela y García, ambas en período de prueba y a cargo de las mismas estudiantes. Se advirtió que quien exponía tendía a confundir los detalles de una causa con los de la otra por su parecido.

\subsection{Caso Sela: acreditar la relación de dependencia oculta}

La demanda original afirmaba que el alimentante trabajaba como portero de un edificio, dato que luego no pudo verificarse en ningún registro. Como prueba se había solicitado que el organismo recaudador (identificado en la reunión como ARCA, denominación actual del organismo fiscal) informara si el demandado estaba inscripto como monotributista. La respuesta oficial indicó que no lo estaba, lo que llevó a inferir que trabaja en relación de dependencia; a partir de ello se solicitó un oficio ampliatorio para que el organismo informe cuánto percibe por esa vía.

En paralelo se libró oficio a la administradora de consorcio para que informe el sueldo que percibe el demandado, oficio que a la fecha de la reunión no había sido respondido. Estos dos oficios (el ampliatorio a ARCA y el dirigido a la administradora) constituyen la prueba pendiente de producción en la demanda.

\subsubsection{El incumplimiento y la retención directa}

Obtenidos alimentos provisorios, el demandado incurrió en incumplimiento. El equipo denunció esa circunstancia y solicitó la retención directa de haberes, oficiando a la administradora para que retenga de su sueldo un monto mensual (cuatrocientos cincuenta pesos como referencia de la cuota); sobre lo ya adeudado, se evaluó solicitar el embargo del sueldo por las cuotas atrasadas.

\begin{definition}{Retención directa y embargo frente al incumplimiento alimentario}
La \textbf{retención directa} opera hacia adelante: consiste en ordenar al empleador o pagador que descuente mensualmente el monto de la cuota vigente y lo deposite en la cuenta de la parte alimentada.

El \textbf{embargo}, en cambio, apunta a las sumas ya adeudadas y no percibidas, y permite afectar haberes o bienes del deudor por el monto de la deuda acumulada.

En este expediente se tramitaron ambas vías en forma simultánea: una para la cuota corriente y otra para lo adeudado.
\end{definition}

Si la administradora finalmente no efectiviza la retención ordenada, ella misma puede quedar alcanzada por responsabilidad solidaria respecto de la deuda alimentaria, en la misma lógica aplicada en otro expediente citado como antecedente por las y los estudiantes.

\subsection{Caso García: el mismo patrón, con Mercado Libre como fuente de prueba}

El caso García repite la lógica del anterior, con una particularidad: en el primer cuatrimestre de 2025 se había omitido diligenciar un oficio dirigido a la plataforma Mercado Libre para determinar si el demandado mantenía fondos en una cuenta de esa plataforma o si había tomado créditos a través de ella, información relevante para establecer su caudal económico. Al no haber respondido la plataforma, se preparó un oficio reiteratorio, que el consultante debía venir a firmar.

También en este expediente hay incumplimiento de alimentos provisorios previos. El demandado manifestó en la causa tener otras dos hijas en la provincia de San Juan y que ya se le retiene el treinta y cinco por ciento de su sueldo, por lo que, según su planteo, no podría afrontar más obligaciones. Una de las estudiantes recordó que en sus anotaciones constaba que ya se había logrado la retención directa y que el expediente había sido remitido a la Defensoría el 7 de agosto, sin que a la fecha de la reunión se conociera el dictamen.

\begin{note}
\textbf{Enseñanza para la práctica profesional.} Reconstruir el patrimonio real de una persona que, en los papeles, no parece tener ingresos es una de las tareas más subestimadas de la práctica de familia. Los oficios a organismos fiscales, administradoras de consorcios y plataformas de comercio electrónico no son un trámite burocrático accesorio, sino la columna vertebral de la prueba en estos expedientes.
\end{note}

% ----------------------------------------------------------
\section{Unidad 4 --- Marcos Calderón: filiación, ejecución de alimentos y extensión a los abuelos}

\subsection{Tres expedientes de una misma familia}

Bajo el nombre del niño beneficiario, Marcos Calderón, conviven tres trámites distintos, lo que llevó a una de las estudiantes a reconocer sentirse ``mareada'' por la cantidad de expedientes correspondientes a una misma persona.

\begin{table}[H]
\centering
\begin{tabularx}{\textwidth}{>{\raggedright\arraybackslash}p{0.27\textwidth} >{\raggedright\arraybackslash}X >{\raggedright\arraybackslash}p{0.2\textwidth}}
\toprule
\textbf{Expediente} & \textbf{Objeto} & \textbf{Estado} \\
\midrule
Filiación & Acreditar la paternidad del alimentante (Mariano) respecto del niño. & Terminado \\
Ejecución de alimentos & Ejecutar la suma que el empleador de Mariano debía haber retenido y no retuvo. & Terminado \\
Alimentos contra los abuelos paternos & Reclamo dirigido en la práctica contra la abuela paterna (Irma). & En preparación \\
\bottomrule
\end{tabularx}
\end{table}

\subsection{Ejecución de alimentos contra el ex empleador}

El segundo expediente se abrió para ejecutar una suma que el empleador de Mariano debía haber retenido de su sueldo en concepto de alimentos y que no retuvo. Se discutió si Mariano continuaba trabajando en el mismo lugar (una heladería) pese a que en algún momento se sostuvo que ya no trabajaba allí. La ejecución avanzó y el ex empleador terminó abonando a la beneficiaria (identificada como Camila) el monto correspondiente.

\begin{important}
\textbf{Responsabilidad del empleador que no retiene los alimentos ordenados judicialmente.} Cuando un juzgado ordena la retención directa de haberes en concepto de alimentos y el empleador, pese a estar notificado, no efectúa esa retención, puede ser compelido a afrontar él mismo el pago de lo que debió retener, tal como ocurrió con el ex empleador de Mariano.
\end{important}

\subsection{Nuevo expediente: alimentos contra la abuela paterna}

El tercer trámite consiste en un reclamo de alimentos dirigido a los abuelos paternos del niño y, en la práctica, contra la abuela (Irma), dado que el abuelo ya no convive con ella. El fundamento fáctico es que Mariano habría dejado de trabajar formalmente, aunque el equipo sostiene que en verdad continúa haciéndolo, sin que hasta el momento se lo hubiera logrado acreditar con fotografías del lugar de trabajo pese a haberlo intentado reiteradamente; por eso no puede seguir ejecutándose contra él.

\textbf{Identificación de la abuela.} Para determinar su domicilio actual se recurrió primero a una dirección que resultó errónea, hasta que una estudiante localizó el perfil de TikTok de la mujer y, a partir de sus publicaciones, pudo identificarse la dirección correcta. Por la misma vía se detectó que la abuela alquila un quincho, lo que indicaría un flujo de ingresos adicional, y que actualmente convive con otra pareja. Al consultar la base de Nosis aparecieron dos personas con el mismo nombre y apellido, lo que generó dudas sobre cuál de los dos números de documento correspondía a la demandada; ante la falta de certeza se decidió avanzar con el de quien, por su nacionalidad paraguaya, resultaba más consistente con el perfil familiar conocido.

\begin{important}
\textbf{Mediación previa obligatoria para demandar alimentos a los abuelos.} A diferencia de los alimentos provisorios, que al tramitarse como medida cautelar no requieren mediación previa, el reclamo de alimentos definitivos contra los abuelos exige agotar la instancia de mediación previa antes de iniciar la demanda. El formulario oficial para solicitarla fue actualizado recientemente (con un nuevo escudo institucional), lo que obligó a rehacer la presentación ya preparada.
\end{important}

Se discutió si, además del reclamo de alimentos futuros contra la abuela, podía intentarse el cobro de la deuda ya generada por el incumplimiento de Mariano. Se aclaró que no: la deuda ya devengada es exclusivamente de Mariano y no puede trasladarse a la abuela, quien solo podría quedar obligada por alimentos futuros, una vez logrado el reclamo en su contra.

\begin{important}
\textbf{Alcance temporal de la responsabilidad de parientes subsidiarios.} Cuando el reclamo alimentario se extiende a un pariente subsidiario (abuelos, hermanos), la obligación de este nace desde que se lo demanda hacia adelante. La deuda ya devengada por el obligado principal y no pagada por él sigue siendo exclusivamente suya y no se traslada retroactivamente al pariente subsidiario.
\end{important}

Por último, en el expediente de filiación también existe una deuda alimentaria de Mariano nunca abonada. Dado que el plazo de prescripción de dos años corre igual, corresponde impulsar allí medidas de ejecución o de embargo en paralelo a la nueva demanda contra la abuela, para evitar que esa deuda prescriba mientras se tramita la mediación.

\begin{note}
\textbf{Enseñanza para la práctica profesional.} Cuando una misma familia genera múltiples expedientes conexos (filiación, ejecución, alimentos contra terceros), es indispensable llevar un registro claro de qué reclamo corresponde a cada causa, qué plazos corren en cada una y cuál es la vía procesal específica, incluida la exigencia o no de mediación previa, aplicable a cada pretensión.
\end{note}

% ----------------------------------------------------------
\section{Cuadro Cornell del Capítulo 1}

\begin{longtable}{>{\raggedright\arraybackslash}p{0.26\textwidth} >{\raggedright\arraybackslash}p{0.64\textwidth}}
\toprule
\textbf{Preguntas / palabras clave} & \textbf{Notas / respuestas} \\
\midrule
\endfirsthead
\toprule
\textbf{Preguntas / palabras clave} & \textbf{Notas / respuestas (continuación)} \\
\midrule
\endhead
\midrule
\multicolumn{2}{r}{\small\itshape continúa en la página siguiente} \\
\endfoot
\bottomrule
\endlastfoot

\textbf{¿Qué es la responsabilidad solidaria entre parientes en materia alimentaria y cuándo se puede invocar?} &
Es la posibilidad de extender el reclamo alimentario a parientes distintos del progenitor obligado en primer término (abuelos, hermanos), en general de manera subsidiaria y ante la imposibilidad de cobrar del obligado principal. Su procedencia por la sola invocación de una situación de vulnerabilidad extrema, sin otro fundamento, resulta incierta según lo relevado en clase, y suele requerir además agotar previamente la vía contra el obligado directo. \\
\addlinespace

\textbf{¿Cuál es el orden de prelación para reclamar alimentos a parientes distintos de los progenitores?} &
Primero se intenta contra los abuelos; solo ante la imposibilidad de dar con ellos corresponde avanzar contra otros parientes colaterales, como el hermano del alimentante (caso Carlovisa). Es una vía excepcional. \\
\addlinespace

\textbf{¿Es imprescindible el CUD para fundar un reclamo alimentario por discapacidad?} &
No. Basta acreditar la discapacidad por otros medios, en línea con el modelo social de la discapacidad, que pone el eje en las barreras que limitan la participación plena de la persona. \\
\addlinespace

\textbf{¿Qué dos extremos debe probar todo reclamo de alimentos?} &
Por un lado, el caudal económico de quien debe pagar (cuánto puede pagar); por otro, la necesidad concreta de quien reclama (cuánto necesita). La dificultad para acreditar el primero, cuando el demandado oculta o disimula sus ingresos, motivó buena parte de la estrategia probatoria de los casos Sela, García y Sobolac. \\
\addlinespace

\textbf{¿Qué diferencia hay entre la retención directa de haberes y el embargo como remedios frente al incumplimiento alimentario?} &
La retención directa opera hacia adelante, ordenando descontar mensualmente la cuota vigente del sueldo del deudor. El embargo apunta a lo ya adeudado y no cobrado, afectando bienes o haberes por el monto de la deuda acumulada. En los casos Sela y García se tramitaron ambas vías de manera simultánea. \\
\addlinespace

\textbf{¿Puede un empleador o administrador que no retiene los alimentos ordenados quedar obligado a pagar?} &
Sí. Si el juzgado ordena la retención directa y el empleador o administrador, estando notificado, no la efectúa, puede quedar obligado a afrontar él mismo el pago de lo que debió retener, como ocurrió con el ex empleador de Mariano en el caso Marcos Calderón. \\
\addlinespace

\textbf{¿Qué requisito previo exige un reclamo de alimentos definitivos contra los abuelos?} &
La mediación previa obligatoria, a diferencia de los alimentos provisorios, que al ser una medida cautelar no la requieren. \\
\addlinespace

\textbf{¿Desde cuándo responde un pariente subsidiario y puede asumir la deuda ya devengada?} &
Su obligación nace desde que se lo demanda, hacia adelante. La deuda ya devengada por el obligado principal sigue siendo exclusivamente de este y no se traslada al pariente subsidiario. \\

\midrule
\multicolumn{2}{>{\raggedright\arraybackslash}p{0.90\textwidth}}{\textbf{Resumen:} Todo reclamo alimentario exige probar el caudal económico del obligado y la necesidad del reclamante. Cuando el obligado principal no puede ser ubicado o no paga, se puede recurrir de manera subsidiaria a parientes (abuelos primero, luego colaterales), con alcance solo hacia adelante y, para definitivos contra abuelos, previa mediación. La prueba del caudal se construye mediante oficios (ARCA, administradoras, plataformas), y frente al incumplimiento se combinan la retención directa (cuota corriente) y el embargo (deuda), con eventual responsabilidad del empleador que no retiene.} \\
\end{longtable}

% ==========================================================
% CAPÍTULO 2
% ==========================================================
\chapter[Incumplimiento, ejecución y coerción]{Incumplimiento, ejecución y coerción del deudor alimentario}

El segundo bloque reúne tres expedientes que comparten un mismo problema de fondo: qué hacer cuando el alimentante incumple de manera sistemática, cuando el conflicto entre las partes contamina el reclamo alimentario o cuando pedir una modificación de la cuota puede terminar perjudicando a quien la pide.

% ----------------------------------------------------------
\section{Unidad 5 --- Ávalos: alimentos trabados por el conflicto en el régimen de comunicación}

\subsection{Antecedentes}

El expediente involucra a Verónica (madre y parte patrocinada), Luciano Sala (padre) y Valentina (hija adolescente, de quince años según se aclaró en clase, aunque en un primer momento se mencionó erróneamente que tenía doce). El equipo tramita en paralelo alimentos provisorios y definitivos a favor de Valentina, ambos actualmente paralizados, y un tercer expediente de régimen de comunicación, muy activo.

\subsection{El nudo del problema: ¿con quién vive Valentina?}

Los alimentos se habían iniciado sobre la base de que Valentina convivía con Verónica. Tras una pelea entre madre e hija, Valentina se mudó a la casa de su padre, quien denunció en el expediente que no correspondía seguir reclamándole alimentos a él si la adolescente ya no vivía con la beneficiaria. Se sucedieron presentaciones cruzadas, cada parte sosteniendo que la joven convivía con el otro progenitor, hasta que se esclareció que Valentina efectivamente reside con Luciano.

En un escrito reciente, Luciano manifestó que no quiere hablar con Verónica, ni verla, ni que ``exista''. El equipo interpretó esta manifestación como una ventaja procesal inesperada: al no tener interés en reclamarle alimentos a Verónica ni en mantener contacto con ella, Luciano no impulsa ningún reclamo alimentario en su contra, lo que resulta favorable para la patrocinada, aunque deja ambos expedientes de alimentos completamente detenidos.

\begin{important}
\textbf{Prescripción de la deuda por alimentos provisorios.} Los alimentos provisorios, al constituir una medida cautelar concedida automáticamente por el juez a pedido de parte, cuentan con una sentencia firme en sí mismos, lo que en principio habilita a reclamar lo adeudado. Sin embargo, el crédito por alimentos prescribe a los dos años según el Código Civil y Comercial, y en este expediente ya había transcurrido aproximadamente un año desde el último impulso, lo que genera urgencia para no perder la posibilidad de cobro.
\end{important}

\subsection{Por qué el equipo decidió no impulsar el cobro}

Pese a existir una vía teórica para reclamar lo adeudado por alimentos provisorios, cada vez que se impulsaba ese reclamo, Luciano respondía con presentaciones en el expediente de régimen de comunicación acusando a Verónica de dañar los sentimientos de Valentina y de ser ``maltratadora'', en lo que se describió como una modalidad de represalia. Según relató la estudiante a cargo, a los dos días de cada avance en el reclamo alimentario aparecía una nueva presentación de la contraparte. Suspender el impulso coincidió con un período de mayor calma en el vínculo entre madre e hija, lo que llevó a mantener ambos expedientes de alimentos en un estado de latencia deliberada.

Ante la pregunta de si corresponde advertirle a la patrocinada que insistir con el reclamo de alimentos puede perjudicar el vínculo con su hija, se aclaró que Verónica ya no podría reclamar alimentos actuales porque Valentina no convive con ella; lo desaconsejable es denunciar el incumplimiento pasado, porque cualquier movimiento en el expediente de alimentos es utilizado por la contraparte (y, según se señaló, lo que dice Luciano suele reaparecer luego en los escritos de la propia adolescente) para reforzar el relato de que Verónica maltrata a Valentina.

\subsection{El expediente de régimen de comunicación: avances, retrocesos y multas}

En una audiencia reciente, celebrada sin la presencia del abogado de Valentina (quien cuenta con patrocinio letrado propio en calidad de adolescente), se alcanzó un acuerdo de régimen de comunicación los días domingo. A los pocos días, Valentina solicitó por escrito ampliar ese régimen a los martes y miércoles, lo que se interpretó en un primer momento como una señal positiva de acercamiento.

Una semana después, Valentina discutió con Verónica (el motivo inmediato habría sido que esta no fue a buscarla al colegio) y revirtió por completo su postura. Resurgió entonces la referencia a un episodio de violencia de género ocurrido en la provincia de Córdoba, ampliamente conocido en ese momento: Verónica le habría dicho a Valentina, como advertencia sobre las consecuencias de seguir viviendo con su padre, algo equivalente a que ``podría terminar como la nena de Córdoba que mataron''. Valentina escribió en el expediente, inmediatamente después del episodio: ``Mamá me trata mal, pero nunca dice nada bueno''. Se comentó la preocupación por el nivel de exposición de la adolescente al conflicto entre sus padres, incluido el hecho de que Valentina interviene en el expediente con una frecuencia que el equipo considera excesiva.

\subsection{La propuesta de terapia vincular}

Frente al estancamiento, el equipo propuso un régimen de terapia vincular en el que madre, padre e hija se sienten frente a un profesional de la psicología para trabajar la dinámica familiar. A la fecha de la reunión ninguna de las tres partes había aceptado la terapia conjunta. El juzgado había requerido a Luciano que informe dónde realizaría una terapia de coparentalidad, lo que quedó sin mayor respuesta concreta, mientras que el equipo debía contestar el traslado correspondiente a la posición de Valentina.

\begin{important}
\textbf{El derecho de niñas, niños y adolescentes a ser oídos.} Se discutió la tensión entre el derecho a ser oídos en los procesos que los afectan, reconocido en la Convención sobre los Derechos del Niño, y el riesgo de que una intervención demasiado frecuente en el expediente genere más malestar que beneficio para la propia adolescente. El equipo concluyó que la vía más adecuada no era pedir que se la excluya del expediente (lo que podría interpretarse como no querer escucharla) sino responder únicamente los planteos formales del padre y no reaccionar a cada intervención espontánea de Valentina.
\end{important}

\subsection{Las multas contra Verónica}

El expediente registra dos sanciones económicas contra Verónica. La primera, una multa ejemplificadora de un millón de pesos, fue impuesta luego de que el juzgado dispusiera que no podía ver a Valentina a solas, sino únicamente con un tercero presente, por los daños emocionales que se le atribuían. El equipo apeló esa multa en el incidente número uno; la apelación se perdió, dado el nivel de malestar manifestado por la propia Valentina.

La segunda sanción, de dos millones de pesos, surge del último dictamen de la Defensoría de Menores, ante la posibilidad de que Verónica vuelva a incurrir en conductas que dañen los sentimientos de la adolescente, fórmula que el equipo calificó de ``muy subjetiva''. El equipo presentó un escrito advirtiendo que, si se llegara a aplicar esa multa, sería apelada, invocando que no pueden apelarse sanciones que todavía no fueron efectivamente ejecutadas.

\begin{note}
\textbf{Enseñanza para la práctica profesional.} En los conflictos de familia de alta litigiosidad, la estrategia procesal no siempre consiste en impulsar todo lo impulsable: a veces la decisión más protectora de los derechos de un niño o adolescente es frenar deliberadamente un reclamo económicamente legítimo, para no alimentar una dinámica de represalias que termina perjudicando al propio beneficiario.
\end{note}

% ----------------------------------------------------------
\section{Unidad 6 --- Sobolac: incumplimiento alimentario crónico, liquidaciones y vía penal}

\subsection{Línea de tiempo del expediente}

Se trata, en palabras de quien expuso, de uno de los expedientes ``más trabados'' del equipo. La consultante (Sobolac) reclama contra su ex pareja (Rossi).

\begin{table}[H]
\centering
\begin{tabularx}{\textwidth}{>{\raggedright\arraybackslash}p{0.2\textwidth} >{\raggedright\arraybackslash}X}
\toprule
\textbf{Fecha} & \textbf{Hito} \\
\midrule
2010 & Separación de hecho de la pareja. \\
2017 & Sentencia de divorcio. \\
2018 & Demanda de alimentos iniciada por la consultante contra Rossi. \\
2021 & Inicio de una causa penal en paralelo. \\
2022 & Primera liquidación aprobada; la causa penal deriva en una probation condicionada a que el demandado se ponga al día con los alimentos. \\
Diciembre de 2022 & Homologación de un acuerdo de pago. \\
Febrero de 2023 & Rossi deja de pagar por completo. \\
Marzo de 2025 & Última liquidación aprobada judicialmente. \\
Julio de 2025 & Nueva liquidación actualizada presentada por el equipo (todo lo devengado durante 2024 y hasta julio de 2025). \\
\bottomrule
\end{tabularx}
\end{table}

% TODO: contenido incompleto o ambiguo en las notas originales (el orden cronológico de la causa penal, la probation y la primera liquidación de 2022 se expuso de manera no lineal)

\begin{important}
\textbf{El acuerdo homologado y su incumplimiento inmediato.} En diciembre de 2022, cuando ya adeudaba aproximadamente un millón cuatrocientos mil pesos entre capital e intereses, Rossi acordó y homologó judicialmente el pago de cuarenta mil pesos mensuales en concepto de alimentos definitivos, más diez mil pesos adicionales para cancelar la deuda acumulada. Cumplió únicamente los pagos de diciembre de 2022 y enero de 2023; desde febrero de ese año dejó de pagar por completo y, según se relató, desapareció.
\end{important}

A partir de ese incumplimiento se sucedieron liquidaciones cada vez más abultadas: la última aprobada judicialmente, en marzo de 2025, ascendía a aproximadamente un millón cuatrocientos mil pesos más doscientos cincuenta mil adicionales por el período reciente, sumas que, se aclaró expresamente, ``hoy no son mucha plata'' en términos de poder adquisitivo. La resolución que siguió a la presentación de julio de 2025 otorgó cinco días a la contraparte para impugnarla o presentar una liquidación alternativa, plazo que probablemente transcurra sin respuesta, lo que llevaría a su aprobación tácita.

Quien expuso resumió la situación así: ``Tenemos liquidaciones, pero nunca se pueden cobrar, porque el tipo está trabajando en algo que no sabemos, tiene sociedades a su nombre y es un fantasma''.

\subsection{La vía penal y los deudores alimentarios morosos}

Ante la reiteración del incumplimiento se evaluó profundizar la vía penal, más allá de la probation ya otorgada. Se hizo referencia a una ley de deudores alimentarios morosos, mencionada de manera aproximada y sin precisión del número exacto, que habilita restricciones como la imposibilidad de renovar la licencia de conducir o de ingresar a determinados espectáculos, además de figurar en un registro público de deudores. Las causas penales no son llevadas directamente por el equipo del patrocinio civil, sino por otra comisión de la Facultad especializada en materia penal, con la que solo se puede interactuar solicitando actualizaciones periódicas, sin poder impulsar activamente los escritos.

% TODO: contenido incompleto o ambiguo en las notas originales (no se precisó el número de la ley de deudores alimentarios morosos)
\begin{note}
Dado que en la exposición no se precisó la norma, a los fines de la práctica profesional se recomienda verificar la ley vigente en la jurisdicción correspondiente antes de invocarla en un escrito.
\end{note}

\subsection{Bienes, sociedades y la nueva pareja}

Rossi tendría sociedades a su nombre, dato que otra estudiante del equipo investigaría en particular, y en un momento se detectaron bienes eventualmente vinculados a su actual pareja (Simonelli). La deuda alimentaria es personal de Rossi, por lo que no sería sencillo trasladar la responsabilidad de pago a su pareja actual por deudas generadas antes de la relación. Se agregó que la hija mayor de edad de la pareja (Marina) se encontraría cursando estudios de ciberseguridad, dato relevante en tanto podría sostener en el futuro un reclamo alimentario propio si continuara estudiando más allá de la mayoría de edad.

\subsection{El nuevo emprendimiento y la lógica del monotributo}

Rossi habría iniciado un emprendimiento de venta de generadores y, según los registros de ARCA, se habría inscripto además para una actividad de venta mayorista de papas fritas desde junio de 2026, bajo la categoría de responsable inscripto (no monotributista). Se razonó que, si una persona se inscribe como responsable inscripto, categoría que exige facturar más y pagar más que el monotributo, sin estar generando ingresos reales, no tendría sentido económico sostener esa carga fiscal; la inferencia del equipo es que, en algún lugar, Rossi efectivamente está facturando y percibiendo ingresos.

Se planteó como línea de acción oficiar nuevamente a ARCA para preguntar puntualmente si el demandado está facturando y, en caso afirmativo, indagar dónde percibe esos fondos, a fin de poder trabar eventualmente un embargo sobre esa fuente de ingresos.

\begin{note}
\textbf{Enseñanza para la práctica profesional.} El caso Sobolac muestra los límites prácticos de la tutela judicial cuando el deudor decide desaparecer del radar formal: contar con una sentencia y liquidaciones aprobadas no equivale a poder cobrar. La estrategia se desplaza hacia la investigación patrimonial (sociedades, inscripciones fiscales, redes sociales) como complemento imprescindible del expediente judicial.
\end{note}

% ----------------------------------------------------------
\section{Unidad 7 --- Segarra: los riesgos de pedir la modificación de la cuota alimentaria}

\subsection{El planteo del consultante}

A diferencia de la mayoría de los expedientes, quien se acercó al patrocinio fue la parte demandada por alimentos (Landolfi), quien manifestó estar pasando considerablemente más tiempo con su hija y consultó si podía modificar el monto que pagaba o directamente dejar de pagarlo.

\subsection{La recomendación: no impulsar el pedido}

El equipo, siguiendo el criterio de la persona a cargo de la supervisión, consideró riesgoso avanzar con una solicitud de modificación de cuota, porque el resultado queda librado por completo a la apreciación judicial. Existe el riesgo cierto de que, en lugar de reducirse, la cuota termine aumentando, sobre todo porque, según se reconoció, el monto que el consultante venía pagando era bajo en relación con sus posibilidades. La indicación al consultante fue que lo pensara, y que si luego de la feria judicial mantenía la misma postura se iniciaría el pedido, pero a riesgo de que se aumentara la cuota.

\begin{important}
\textbf{El riesgo procesal de pedir una modificación de cuota.} Solicitar la modificación de una cuota alimentaria abre nuevamente a discusión todo el cuadro de necesidades del beneficiario y de capacidad económica del alimentante, sin que el resultado esté garantizado en el sentido pretendido por quien lo solicita. Un pedido de reducción puede derivar, luego de la valoración judicial de la prueba, en un aumento de la cuota vigente, por lo que la recomendación profesional debe incluir siempre una advertencia clara sobre ese riesgo antes de avanzar.
\end{important}

\subsection{Desenlace}

El consultante nunca volvió a responder a los llamados del equipo (se mencionaron alrededor de cuarenta llamadas perdidas), por lo que el expediente quedó, en los hechos, sin impulso.

\begin{note}
\textbf{Enseñanza para la práctica profesional.} Explicar con claridad a la persona consultante los riesgos de una acción legal, incluso cuando ello implica desalentar el pedido que trajo originalmente al patrocinio, forma parte del deber de información y buena fe hacia quien confía su caso al abogado o abogada.
\end{note}

% ----------------------------------------------------------
\section{Cuadro Cornell del Capítulo 2}

\begin{longtable}{>{\raggedright\arraybackslash}p{0.26\textwidth} >{\raggedright\arraybackslash}p{0.64\textwidth}}
\toprule
\textbf{Preguntas / palabras clave} & \textbf{Notas / respuestas} \\
\midrule
\endfirsthead
\toprule
\textbf{Preguntas / palabras clave} & \textbf{Notas / respuestas (continuación)} \\
\midrule
\endhead
\midrule
\multicolumn{2}{r}{\small\itshape continúa en la página siguiente} \\
\endfoot
\bottomrule
\endlastfoot

\textbf{¿Cuál es el plazo de prescripción de la deuda alimentaria según el Código Civil y Comercial?} &
Dos años. Ese plazo resultó determinante en más de un expediente (Ávalos, Sobolac) para decidir con urgencia si convenía o no impulsar el cobro de sumas adeudadas antes de que operara la prescripción. \\
\addlinespace

\textbf{¿Por qué, en algunos casos, el equipo decidió deliberadamente no impulsar un reclamo alimentario pese a tener derecho a hacerlo?} &
Porque impulsar el reclamo generaba, en la práctica, una represalia inmediata de la contraparte en un expediente conexo (régimen de comunicación), perjudicando el vínculo entre madre e hija más de lo que beneficiaba económicamente a la parte patrocinada, como se explicó extensamente en el caso Ávalos. \\
\addlinespace

\textbf{¿Cómo se resolvió la tensión entre el derecho del adolescente a ser oído y su sobreexposición en el expediente?} &
No se pidió excluir a la adolescente del expediente, lo que podría interpretarse como no querer escucharla; se optó por responder únicamente los planteos formales del padre y no reaccionar a cada intervención espontánea de Valentina. \\
\addlinespace

\textbf{¿Qué alternativas quedan cuando las liquidaciones aprobadas no se pueden cobrar?} &
Se desplaza la estrategia hacia la investigación patrimonial (sociedades, inscripciones fiscales, redes sociales) y se pueden requerir medidas de coerción indirecta o profundizar la vía penal. En Sobolac se planteó oficiar nuevamente a ARCA para saber si el demandado factura y dónde percibe los fondos, a fin de trabar un embargo. \\
\addlinespace

\textbf{¿Qué restricciones asociaba la clase a los deudores alimentarios morosos?} &
Imposibilidad de renovar la licencia de conducir, de acceder a determinados espectáculos y figuración en un registro público de deudores. No se precisó el número de la norma, por lo que debe verificarse antes de invocarla. \\
\addlinespace

\textbf{¿Quién responde por la deuda alimentaria anterior a la nueva pareja del deudor?} &
La deuda alimentaria es personal del deudor (Rossi); no sería sencillo trasladar la responsabilidad a su pareja actual por deudas generadas antes de la relación. \\
\addlinespace

\textbf{¿Qué riesgo implica solicitar la modificación de una cuota alimentaria ya fijada?} &
El riesgo de que, al reabrirse a discusión todo el cuadro de necesidades y capacidad económica, el resultado no sea el buscado: una petición de reducción puede derivar en un aumento de la cuota, como se advirtió en el caso Segarra. La recomendación profesional debe incluir siempre esa advertencia. \\

\midrule
\multicolumn{2}{>{\raggedright\arraybackslash}p{0.90\textwidth}}{\textbf{Resumen:} La deuda alimentaria prescribe a los dos años, lo que obliga a decidir con urgencia si impulsar el cobro. En conflictos de alta litigiosidad puede ser más protector frenar deliberadamente un reclamo legítimo para evitar represalias que perjudiquen al beneficiario. Frente al deudor que desaparece del radar formal, la estrategia se apoya en la investigación patrimonial, y antes de pedir una modificación de cuota debe advertirse al consultante que el resultado puede ser un aumento.} \\
\end{longtable}

% ==========================================================
% CAPÍTULO 3
% ==========================================================
\chapter[Comunicación, violencia familiar y protección]{Régimen de comunicación, violencia familiar y protección de niñas, niños y adolescentes}

El tercer bloque reúne expedientes en los que el eje ya no es principalmente económico, sino vincular: cómo ordenar la convivencia y el contacto entre adultos y niños cuando el vínculo de pareja se rompió, y cómo actuar cuando existen antecedentes de violencia o dudas sobre la salud mental de alguno de los progenitores.

% ----------------------------------------------------------
\section{Unidad 8 --- Walter: régimen de comunicación entre convivientes y desocupación del inmueble}

\subsection{Antecedentes}

Walter se acercó a la Facultad para establecer un régimen de comunicación con su ex pareja, con quien continuaba conviviendo bajo el mismo techo. Según relató, ella amenazaba de manera reiterada con irse de la casa junto con el hijo en común, lo que generaba un clima de conflicto permanente.

\subsection{El primer acuerdo}

El equipo elaboró un acuerdo de convivencia que fue firmado por ambas partes (debiendo intervenir también, no sin dificultades, la abogada de la ex pareja de Walter), con pautas de convivencia en la casa y previsiones sobre cómo proceder ante eventuales viajes, entre otros puntos. El acuerdo contemplaba además que Walter aguardaría un tiempo razonable para que su ex pareja consiguiera un lugar propio donde vivir.

A los pocos días de firmado, Walter volvió a consultar cómo exigir que su ex pareja se retirara de la vivienda. El equipo le explicó que, dado lo reciente del acuerdo y el compromiso asumido de otorgarle un plazo razonable, no correspondía exigir el desalojo de manera inmediata.

\begin{important}
\textbf{Opciones evaluadas frente al conflicto de convivencia.} Antes de recurrir a una acción judicial de desalojo, considerada desproporcionada dado que el hijo en común convive con ambos en la misma vivienda, se barajaron alternativas. Entre ellas, el envío de una carta documento, bajo la hipótesis de que una comunicación formal de ese tipo podría ser suficiente para motivar la mudanza de la ex pareja sin necesidad de un proceso judicial de desalojo.
\end{important}

\subsection{El giro: el viaje a España}

Tiempo después, Walter informó que había cambiado de parecer: proyectaba viajar a España a fin de año junto con su hijo, con el consentimiento de su ex pareja, quien permanecería viviendo en la casa de Walter. Se discutió si correspondía cobrarle algún alquiler por esa ocupación, sin arribar a una definición final, dado que, según se interpretó, Walter estaría dispuesto incluso a no cobrar alquiler con tal de garantizar que ella no obstaculice el viaje con su hijo.

Un compañero observó que, de tratarse de un terreno con varias viviendas (al parecer ocupado también por otros integrantes de la familia), cualquier autorización de viaje que se firme ahora podría ser utilizada luego por la ex pareja o su abogado para reclamar, en el futuro, un contrato de comodato de larga duración sobre el inmueble.

\begin{important}
\textbf{Autorización de viaje y centro de vida del niño.} Cuando un progenitor proyecta radicarse en el exterior con un hijo o hija, corresponde formalizar tanto el consentimiento del otro progenitor para el viaje como, en su caso, un nuevo acuerdo o trámite que fije el centro de vida del niño junto con quien viaja. El equipo previó convocar a Walter en diciembre, antes del receso de la Facultad, para tramitar con anticipación el permiso de viaje y cualquier otro trámite legal necesario, de modo de evitar imprevistos antes de la partida.
\end{important}

\begin{note}
\textbf{Enseñanza para la práctica profesional.} Un mismo conflicto de convivencia puede transformarse por completo en pocas semanas: de una tensión doméstica que amenazaba con derivar en un desalojo, a un escenario de cooperación forzada motivado por un interés compartido, el viaje del hijo. La tarea del abogado o abogada de familia no es solo redactar el acuerdo que corresponde a la situación de hoy, sino anticipar cómo ese mismo acuerdo puede ser usado o reinterpretado si la situación cambia.
\end{note}

% ----------------------------------------------------------
\section{Unidad 9 --- Elías: violencia de género, medidas de restricción y régimen de comunicación}

\subsection{Dos expedientes conexos}

El consultante, Elías, ingresó al patrocinio en 2024 con dos expedientes: uno de violencia, iniciado por su ex pareja (Soria) mediante una denuncia ante la Oficina de Violencia Doméstica (OVD), y otro de régimen de comunicación con el niño en común, promovido posteriormente por el propio Elías.

\subsection{El expediente de violencia: medidas sucesivas}

A lo largo de 2024, en el marco de la denuncia de Soria, se dictaron sucesivas medidas de restricción respecto de Elías, tanto en relación con ella como con el niño. El equipo decidió no impulsar activamente ese expediente, con el objetivo de que las medidas de restricción no fueran renovadas, y concentrar los esfuerzos en el nuevo expediente de régimen de comunicación, dado que Elías no veía a su hijo desde 2024.

\begin{important}
\textbf{Condiciones impuestas judicialmente a Elías.} El juzgado le impuso tres condiciones: (i) realizar un tratamiento por violencia; (ii) respetar las medidas de restricción vigentes; y (iii) realizar un tratamiento psicológico o psiquiátrico en un establecimiento especializado, en este caso el hospital Borda. Elías cumplió las dos primeras, pero nunca llegó a concretar la tercera, en parte porque el hospital tampoco respondió a las gestiones iniciadas por el propio equipo del patrocinio.
\end{important}

\subsection{El expediente de régimen de comunicación}

Al vencer las medidas de restricción hacia abril de 2025 sin haber sido renovadas, el equipo promovió la demanda de régimen de comunicación. El principal obstáculo procesal fue la notificación: la primera cédula resultó negativa y, al indagar el motivo, se descubrió que la ex pareja de Elías no tiene, como se había asumido al presentar la demanda, domicilio en la Ciudad Autónoma de Buenos Aires, sino en la provincia de Buenos Aires. El propio expositor reconoció el error como propio por no haber verificado ese dato con precisión.

Se evaluó notificar a través del domicilio electrónico de las abogadas que representan a la contraparte en el expediente conexo de violencia, mediante un pedido de vinculación entre ambas causas, aunque se discutió si esa vía era jurídicamente correcta dado que la notificación de una demanda debe ser, en principio, personal.

\begin{important}
\textbf{Vinculación de expedientes y domicilio electrónico.} Cuando existen dos expedientes conexos entre las mismas partes, puede solicitarse su vinculación para aprovechar, entre otras cosas, el domicilio electrónico ya constituido por los letrados intervinientes en uno de ellos a los fines de las notificaciones en el otro. Esta vía no reemplaza necesariamente el requisito de notificación personal cuando la ley así lo exige, y las abogadas de la contraparte en el expediente conexo podrían válidamente rechazar la vinculación.
\end{important}

\subsection{El objetivo perseguido: una videollamada}

Una vez lograda la notificación, el pedido concreto es autorizar una comunicación por videollamada entre Elías y su hijo, algo que el equipo consideraba razonablemente alcanzable si se lograba notificar a la contraparte.

\subsection{El posible abandono del caso}

Sobre el cierre de la exposición se informó que Elías había manifestado, mediante un mensaje de WhatsApp puesto a disposición del equipo, que había dejado de trabajar y que ya no quería seguir pagando alimentos, razón por la cual no se sabía con certeza si continuaría impulsando el caso. Ese mensaje ya había sido comentado con la persona a cargo de la supervisión, quien indicó que, si Elías no vuelve a presentarse, el expediente debe dejarse tal como está.

\begin{note}
\textbf{Enseñanza para la práctica profesional.} El caso combina la violencia de género, con sus medidas de protección y condiciones de tratamiento, y el derecho del progenitor no conviviente a mantener contacto con su hijo. La responsabilidad profesional exige sostener ambos frentes con cuidado, verificando con rigor cuestiones básicas pero decisivas, como el domicilio real de la contraparte, antes de iniciar cualquier trámite de notificación.
\end{note}

% ----------------------------------------------------------
\section{Unidad 10 --- Díaz: salud mental parental, informes interdisciplinarios y protección del niño}

Este expediente se presentó como uno de los más complejos de la reunión, con la advertencia previa de que insumiría más tiempo que los anteriores.

\subsection{Antecedentes}

La consultante, Díaz, había iniciado en 2020 una denuncia en la Oficina de Violencia Doméstica que no llegó a prosperar. Posteriormente su ex pareja (Gaspari, padre de la niña en común) inició una nueva denuncia en la misma oficina, esta vez contra Díaz. El expediente relevante para el patrocinio se inicia en 2023 y, a lo largo de su tramitación, Díaz tuvo dos abogados distintos, ambos de los cuales renunciaron a la representación por desavenencias con su patrocinada, vinculadas, según se relató, a que cuando el juzgado resuelve algo en su contra suele responsabilizar al propio abogado.

Se aclaró que la denuncia por violencia iniciada por Díaz en 2020 sí generó un expediente formal, aunque el equipo no tiene acceso a su contenido; distinto es el caso de la denuncia presentada por Gaspari, que avanzó hasta derivar en el incidente que restringió el contacto de Díaz con su hija.

\subsection{Cronología de intimaciones y evaluaciones}

Entre 2023 y 2025 se sucedieron intimaciones a Díaz para que se presentara, se notificara y aceptara un régimen de comunicación con el padre de la niña, además de requerimientos para que iniciara un tratamiento psicológico, que ella se negaba sistemáticamente a realizar. Ante la dificultad para notificarla, el juzgado optó por notificarla directamente por correo electrónico. En marzo de 2026 se le otorgó un nuevo plazo prudencial para presentarse, luego de que invocara reiteradamente no estar trabajando y no poder hacerlo.

\subsection{Los informes que fueron delineando el diagnóstico}

El equipo repasó, en el orden en que se incorporaron al expediente, los informes técnicos que configuraron la evaluación desfavorable respecto de la capacidad de Díaz para tener a su hija a cargo:

\begin{itemize}
    \item \textbf{Informe del Ministerio Público Tutelar sobre el padre.} Evaluó al padre y concluyó que no presentaba signos psicopáticos ni alteración de la realidad, recomendando una revinculación paulatina entre la niña y su madre.
    \item \textbf{El episodio de las manchas en el rostro.} Una trabajadora social del hospital Vélez Sarsfield se comunicó con el colegio de la niña luego de que Díaz insistiera, en una consulta médica, en que su hija presentaba manchas en el rostro provocadas en el colegio, pese a que el personal médico no las constató. Díaz sostenía además que la niña le había sido entregada al padre durante el horario escolar sin su consentimiento y que existía una suerte de connivencia entre las psicólogas, la directora y las maestras del establecimiento en contra suya.
    \item \textbf{Informe del colegio.} Informó que había convocado a Díaz en reiteradas oportunidades y que ella se negaba sistemáticamente a asistir, y que la niña acumulaba un porcentaje elevado de inasistencias (mencionado en clase en el orden del 70\,\%).
    \item \textbf{Dictamen de la Defensoría de Menores.} Describió una escasa red de apoyo familiar, ausentismo escolar, desajuste discursivo y conductual de la madre y obstrucción a la intervención institucional, además de señalar que Díaz involucra a la niña en los conflictos con el padre. Observó también que se negaba a ser evaluada psicológicamente, hasta que finalmente se recomendó una evaluación interdisciplinaria de su salud mental.
    \item \textbf{Consejo de Niñas, Niños y Adolescentes zonal.} La Defensoría zonal recomendó mantener la orden de alejamiento respecto de Díaz, extendida no solo al padre y a la niña sino también a la abuela paterna, con quien el padre y la niña conviven, dado que se relevaron episodios en los que Díaz habría manifestado, según consta en el expediente, ``tintes persecutorios''.
    \item \textbf{La aparición de la hermana mayor.} En determinado momento apareció en el expediente la hija mayor de Díaz, de una relación anterior, de treinta años, con patrocinio letrado propio, manifestando no tener inconveniente en ser evaluada. El padre denunció que esta hija le enviaba de manera cotidiana decenas de mensajes (alrededor de ciento cincuenta capturas de pantalla) insistiendo en ver a la niña, lo que motivó un pedido para que también se abstuviera de acercarse hasta contar con su propia evaluación psicológica programada.
    \item \textbf{Último informe psicológico de la niña.} Señaló que, desde que convive con su padre, se la observa considerablemente mejor, y que su comportamiento cambia de forma notoria cuando concurre acompañada por su madre, respecto de cuando se presenta con su padre.
\end{itemize}

\begin{important}
\textbf{Un diagnóstico construido por convergencia de informes, no por un único peritaje.} Frente a la pregunta de si existía una evaluación psicológica formal que calificara a Díaz de un modo determinado, el equipo aclaró que ningún informe la describe en términos tan categóricos, pero que la convergencia entre el informe del colegio, el de la Defensoría zonal, el del Ministerio Público Tutelar y las evaluaciones psicológicas del padre y de la niña señalan, en conjunto, rasgos persecutorios y paranoides y una distorsión de la percepción de la realidad, ejemplificada en el episodio de las manchas inexistentes.
\end{important}

\subsection{La postura de la consultante y el pronóstico del equipo}

En la entrevista personal con el equipo, Díaz manifestó que no está ``loca'' y que no concurrirá a ningún hospital para ser evaluada. Frente a ese nivel de resistencia, el equipo coincidió en que la única estrategia con alguna probabilidad de éxito consiste en convencerla de iniciar un tratamiento terapéutico y demostrar progresivamente en el expediente los avances de esa terapia, como condición previa e indispensable para plantear más adelante una revinculación paulatina como la que ya se dispuso en otros expedientes similares. El pronóstico expresado fue que, con todos estos antecedentes, ``no tiene mucha defensa''.

Se anticipó además una dificultad adicional: el juzgado ya había ordenado a Díaz entregar determinados objetos personales de la niña (incluyendo, según se mencionó, una mascota), orden que ella manifestó no estar dispuesta a cumplir, lo que revela una dificultad más profunda para comprender el carácter obligatorio de una orden judicial.

\begin{note}
\textbf{Enseñanza para la práctica profesional.} El caso confronta a quienes ejercen la abogacía de familia con uno de los dilemas más delicados de la profesión: patrocinar a una persona cuya narrativa de los hechos puede estar distorsionada, sin que ello implique abandonar su derecho de defensa, pero exigiendo señalarle con honestidad cuáles son las condiciones reales para poder avanzar en su reclamo. La multiplicidad de informes técnicos convergentes ilustra cómo el expediente de familia se construye muchas veces a partir de un entramado interdisciplinario (escuela, hospital, defensorías, peritos) y no de una sola prueba decisiva.
\end{note}

% ----------------------------------------------------------
\section{Cuadro Cornell del Capítulo 3}

\begin{longtable}{>{\raggedright\arraybackslash}p{0.26\textwidth} >{\raggedright\arraybackslash}p{0.64\textwidth}}
\toprule
\textbf{Preguntas / palabras clave} & \textbf{Notas / respuestas} \\
\midrule
\endfirsthead
\toprule
\textbf{Preguntas / palabras clave} & \textbf{Notas / respuestas (continuación)} \\
\midrule
\endhead
\midrule
\multicolumn{2}{r}{\small\itshape continúa en la página siguiente} \\
\endfoot
\bottomrule
\endlastfoot

\textbf{¿Cómo se ordena una convivencia conflictiva sin llegar al desalojo?} &
Mediante un acuerdo de convivencia firmado por ambas partes, con pautas y previsiones sobre viajes, que otorga un plazo razonable para que la ex pareja consiga otro lugar. Se consideró desproporcionada la acción de desalojo por convivir el hijo con ambos, y se barajó el envío de una carta documento (caso Walter). \\
\addlinespace

\textbf{¿Qué debe formalizarse cuando un progenitor proyecta radicarse en el exterior con un hijo?} &
El consentimiento del otro progenitor para el viaje y, en su caso, un nuevo acuerdo o trámite que fije el centro de vida del niño junto con quien viaja, con la debida anticipación. Debe preverse además que lo firmado puede ser reinterpretado en el futuro (por ejemplo, para reclamar un comodato de larga duración). \\
\addlinespace

\textbf{¿Qué condiciones debe cumplir el agresor antes de reanudar el contacto?} &
En el caso Elías, el juzgado impuso tres condiciones: tratamiento por violencia, respeto de las medidas de restricción vigentes y tratamiento psicológico o psiquiátrico en un establecimiento especializado (hospital Borda). Cumplió las dos primeras y nunca concretó la tercera. \\
\addlinespace

\textbf{¿Para qué sirve vincular dos expedientes conexos y qué límite tiene?} &
Permite aprovechar el domicilio electrónico de los letrados constituido en uno de ellos para notificar en el otro. No reemplaza necesariamente la notificación personal de la demanda cuando la ley la exige, y la contraparte podría rechazar la vinculación. \\
\addlinespace

\textbf{¿Qué rol cumple el Ministerio Público Tutelar o la Defensoría de Menores e Incapaces en los expedientes de familia?} &
Interviene emitiendo dictámenes técnicos sobre la situación de niñas, niños y adolescentes o de personas con capacidad restringida, evaluando a las partes y recomendando medidas de protección, revinculación o restricción. En los casos Ávalos y Díaz sus dictámenes resultaron determinantes tanto para la imposición de sanciones como para la construcción del cuadro probatorio sobre la aptitud parental. \\
\addlinespace

\textbf{¿Cómo pondera la justicia múltiples informes sobre un progenitor?} &
Por convergencia: aunque ningún informe la calificaba en términos categóricos, el conjunto (colegio, Defensoría zonal, Ministerio Público Tutelar y evaluaciones del padre y de la niña) señalaba rasgos persecutorios y paranoides y una distorsión de la percepción de la realidad (caso Díaz). \\
\addlinespace

\textbf{¿Qué estrategia tiene probabilidad de éxito frente a una progenitora que rechaza toda evaluación?} &
Convencerla de iniciar un tratamiento terapéutico y demostrar progresivamente en el expediente sus avances, como condición previa para plantear una revinculación paulatina. \\

\midrule
\multicolumn{2}{>{\raggedright\arraybackslash}p{0.90\textwidth}}{\textbf{Resumen:} En los conflictos vinculares el eje deja de ser económico y pasa a ser cómo ordenar la convivencia y el contacto con los hijos. Los acuerdos deben redactarse anticipando su reinterpretación futura; en la violencia familiar, el contacto se condiciona al cumplimiento de tratamientos y medidas; y la situación parental se evalúa por la convergencia de informes interdisciplinarios, con intervención determinante del Ministerio Público Tutelar y las defensorías. Verificar datos básicos, como el domicilio real, evita errores de notificación.} \\
\end{longtable}
\end{document}
